% Source-derived chapter 02: Background on parabolic bundles
% Original source: geometric-local-systems-general-curves-isomonodromy
\chapter{Background on parabolic bundles}
\label{geometric-local-systems-general-curves-isomonodromy-chapter-02}
\source{geometric-local-systems-general-curves-isomonodromy}

\begin{remark}[Source section]
\label{geometric-local-systems-general-curves-isomonodromy-chapter-02-source}
\source{geometric-local-systems-general-curves-isomonodromy}
\source{geometric-local-systems-general-curves-isomonodromy}
This chapter reproduces the corresponding section of the retrieved
LaTeX source, including its definitions, statements, examples, and
proofs. It has not yet been translated into Lean.
\notready
\end{remark}

% The source section title was: Background on parabolic bundles
\label{section:parabolic}
\source{geometric-local-systems-general-curves-isomonodromy}

We now review some basics on parabolic sheaves and parabolic bundles, primarily following the notation of
\cite[\S2.3]{BHH:parabolic}.
Some useful references include 
\cite[Part 3, \S1]{seshadri:fibres-vectoriels-sur-les-courbes-algebriques},
\cite[\S1 and \S3]{yokogawa:infinitesimal-deformation}, and
\cite[\S2]{bodenY:moduli-spaces-of-parabolic-higgs-bundles}.
Let $C$ be a smooth proper curve and $D\subset C$ a reduced divisor. Loosely speaking, a parabolic bundle is a vector bundle on $C$ together with an additional
filtration of the fibers over points of the given divisor, weighted by an increasing sequence of real numbers in $[0,1)$.

\subsection{Definition of parabolic bundles}
\label{subsubsection:definition-parabolic}
\source{geometric-local-systems-general-curves-isomonodromy}

\begin{definition}
\source{geometric-local-systems-general-curves-isomonodromy}
\notready
	\label{definition:parabolic-bundle}
\source{geometric-local-systems-general-curves-isomonodromy}
	Let $E$ be a vector bundle over a curve $C$.
	Let $D = x_1 + \cdots + x_n \subset C$ denote a divisor, with the $x_i$ distinct.
	A {\em quasiparabolic structure} on $E$ over $D$ is a strictly decreasing
	filtration of subspaces
	\begin{align*}
		E_{x_j} = E_j^1 \supsetneq E_j^2 \supsetneq \cdots \supsetneq E_j^{n_j+1}
		= 0
	\end{align*}
	for each $1 \leq j \leq n$.
	A {\em parabolic structure} on $E$ over $D$ is a quasiparabolic
	structure together with
	$n$ sequences of real numbers
	\begin{align*}
		0 \leq \alpha^1_j < \alpha^2_j \cdots < \alpha^{n_j}_j < 1
	\end{align*}
	for $1 \leq j \leq n$.
	A {\em parabolic bundle} is a vector bundle with a parabolic structure.
	The collection $\{\alpha^i_j\}$ are called the {\em weights} of the
	parabolic bundle. We say the parabolic bundle has {\em rational weights} if
	all $\alpha^i_j$ are rational numbers.
	We often notate the data of a parabolic bundle simply as $E_\star$
	instead of $(E, \{E^i_j\}, \{\alpha^i_j\})$.
\end{definition}

\begin{definition}[Parabolic degree and slope]
\source{geometric-local-systems-general-curves-isomonodromy}
\notready
	\label{definition:}
\source{geometric-local-systems-general-curves-isomonodromy}
	Let $E_\star := (E, \{E^i_j\}, \{\alpha^i_j\})$. We define the {\em
	parabolic degree} of $E_\star$ as
\begin{align*}
	\on{par-deg}(E_\star) := \deg(E) + \sum_{j=1}^n \sum_{i=1}^{n_j}
	\alpha^i_j \dim(E_j^i/E_j^{i+1}).
\end{align*}
We define the {\em parabolic slope} as $\mu_\star(E_\star) :=
\on{par-deg}(E_\star)/\rk(E_\star)$.
\end{definition}

\subsection{Definition of parabolic sheaves}
\label{subsubsection:parabolic-sheaf}
\source{geometric-local-systems-general-curves-isomonodromy}

In order to later apply Serre duality, we will not only need parabolic bundles, but also so-called
coparabolic bundles. Coparabolic bundles are examples of parabolic sheaves,
which we define next.
The definitions in this subsection follow
\cite{yokogawa:infinitesimal-deformation} and
\cite{bodenY:moduli-spaces-of-parabolic-higgs-bundles}.

\begin{definition}
\source{geometric-local-systems-general-curves-isomonodromy}
\notready
	\label{definition:filtered-O_X-module}
\source{geometric-local-systems-general-curves-isomonodromy}
	Let $X$ be a scheme and $D \subset X$ an effective Cartier divisor.
	Let $\mathbb R$ denote the category whose objects are real numbers with
	a single morphism $i^{\alpha,\beta}:\alpha\to\beta$ if $\alpha \geq \beta$ and no
	morphisms otherwise.
	Let $\mathscr M_X$ denote the category of sheaves of $\mathscr{O}_X$-modules.
A {\em $\mathbb R$-filtered $\mathscr O_X$-module} is a functor $E: \mathbb R \to \mathscr M_X$.
Notationally, we use $E_\star$ to denote the functor $E$, so $E_\alpha := E(\alpha)$.
We write $i_E^{\alpha,\beta} := E(i^{\alpha,\beta})$.
\end{definition}
\begin{example}
\source{geometric-local-systems-general-curves-isomonodromy}
\notready
	\label{example:}
\source{geometric-local-systems-general-curves-isomonodromy}
	Define $E[\alpha]_\star$ as the $\mathbb{R}$-filtered $\mathscr O_X$-module given by
$E[\alpha]_\beta := E_{\alpha + \beta}$
and $i^{\beta,\gamma}_{E[\alpha]} := i_E^{\beta + \alpha, \gamma + \alpha}$.
Let $i^{[\alpha,\beta]}_E : E[\alpha]_\star \to E[\beta]_\star$ denote the
natural transformation whose value on $\gamma$ is $i_E^{\alpha + \gamma, \beta +
\gamma}$.
For $f: E_\star \to F_\star$, we use $f[\alpha] : E[\alpha]_\star \to F[\alpha]_\star$ for
the natural induced map.
By abuse of notation, we will frequently write $E$ in place of $E_0$.
\end{example}
\begin{definition}
\source{geometric-local-systems-general-curves-isomonodromy}
\notready
	\label{definition:parabolic-sheaf}
\source{geometric-local-systems-general-curves-isomonodromy}
A {\em parabolic sheaf} on $X$ with respect to $D$ is
a $\mathbb R$-filtered $\mathscr O_X$-module $E_\star$ equipped with an isomorphism
\begin{align*}
	j_E : E_\star \otimes \mathscr O_X(-D) \overset{\sim}{\to} E[1]_\star
\end{align*}
 such that $i_E^{[1,0]} \circ j_E = \on{id}_{E_\star}
\otimes i_D: E_\star \otimes \mathscr O_X(-D) \to E_\star$, where $i_D: \mathscr{O}_X(-D)\to \mathscr{O}_X$ is the natural inclusion.

A natural transformation of parabolic $\mathscr O_X$-modules $f: E_\star \to
F_\star$ is a {\em parabolic morphism} if
\begin{equation}
	\label{equation:}
\source{geometric-local-systems-general-curves-isomonodromy}
	\begin{tikzcd} 
		E_\star \otimes \mathscr O_X(-D) \ar {r}{f\otimes
		\on{id}} \ar {d}{j_E} & F_\star \otimes \mathscr
		O_X(-D) \ar {d}{j_F} \\
		E[1]_\star \ar {r}{f[1]} & F[1]_\star
\end{tikzcd}\end{equation}
commutes.

Let $\on{Hom}(E_\star, F_\star)$ denote the set of parabolic morphisms.
We let $\sheafhom(E_\star, F_\star)$ denote the sheaf of homomorphisms 
defined by taking $\sheafhom(E_\star, F_\star)(U):= \on{Hom}(E_\star|_U,
F_\star|_U).$
We also define a parabolic sheaf
$\sheafhom(E_\star, F_\star)_\star$ by taking
$\sheafhom(E_\star, F_\star)_\alpha := \sheafhom(E_\star, F[\alpha]_\star)$
and the $i_{\sheafhom(E_\star, F_\star)}^{\alpha,\beta}$ for $\alpha \geq \beta$
to be the natural maps
induced by $i_F^{[\alpha, \beta]}: F[\alpha]_\star \to F[\beta]_\star$.

We define $\on{End}(E_\star) := \on{Hom}(E_\star, E_\star)$
and use $\sheafend(E_\star) := \sheafhom(E_\star, E_\star)$.
\end{definition}



\begin{example}
\source{geometric-local-systems-general-curves-isomonodromy}
\notready
	\label{example:parabolic}
\source{geometric-local-systems-general-curves-isomonodromy}
	Any parabolic vector bundle $(E, \{E^i_j\}, \{\alpha^i_j\})$ defines a
	parabolic sheaf as follows. For $0 \leq \alpha < 1$,
	define
	\begin{align*}
	E_\alpha := \cap_{j=1}^n \ker(E \to
	E_{x_j}/E^{\beta(\alpha, j)}_j).
	\end{align*}
	where $\beta(\alpha, j)=\min(i: \alpha^i_j\geq \alpha),$ for $\alpha\leq \max_i(\alpha^i_j)$ and taking $\beta(\alpha, j)=n_j+1$ for $1>\alpha>\max_i(\alpha^i_j)$.
	For $n\in \mathbb{Z}$ set $E_{x+n} := E_x(-nD)$ and take $i^{\alpha,\beta}_E$ as the
	natural inclusions. We will refer to parabolic vector bundles and their associated parabolic sheaves interchangeably.
\end{example}

\begin{remark}
\source{geometric-local-systems-general-curves-isomonodromy}
\notready
	\label{remark:end-definition}
\source{geometric-local-systems-general-curves-isomonodromy}
	It follows from the definition of $\on{End}(E_\star)$ that 
	$\sheafend(E_\star) \subset \sheafend(E)$ is the coherent subsheaf
corresponding to those endomorphisms preserving the quasiparabolic filtration
$\{E^i_j\}$ at each point $x_j$ of $D$.
\end{remark}

\begin{example}
\source{geometric-local-systems-general-curves-isomonodromy}
\notready
	\label{example:vector-bundle}
\source{geometric-local-systems-general-curves-isomonodromy}
	Every vector bundle $E$ defines a parabolic bundle on $X$ with respect
	to $D$ by taking the quasiparabolic structure
	$E_{x_j} = E_j^1 \subset E_{j}^2 = 0$ with $\alpha_j^1 = 0$.
	In turn, this defines a parabolic sheaf by \autoref{example:parabolic}.
	We say such a parabolic bundle has trivial parabolic structure.
\end{example}
\begin{remark}
\source{geometric-local-systems-general-curves-isomonodromy}
\notready
	\label{remark:map-vector-bundle-to-parabolic}
\source{geometric-local-systems-general-curves-isomonodromy}
	If $V$ is a vector bundle on a scheme $X$ and $E_\star$ is a parabolic
	sheaf, we write a map $V \to E_\star$ to denote a map $V_\star \to
	E_\star$, where $V_\star$ is the parabolic sheaf corresponding to $V$ as
	in \autoref{example:vector-bundle}.
\end{remark}


In addition to parabolic vector bundles, we will also require coparabolic vector
bundles, in order to use Serre duality.
The essential idea is that while the parabolic sheaves associated to parabolic vector bundles are unchanged in intervals of
the form $(\alpha^i, \alpha^{i+1}]$, coparabolic vector bundles are unchanged in
intervals of the form $[\alpha^i, \alpha^{i+1})$.
I.e., parabolic vector bundles can be viewed as lower semicontinuous functors, taking
the discrete topology on the set of isomorphism classes of vector bundles, while coparabolic vector
bundles 
can be viewed as upper semicontinuous functors, see \cite[Figure
1]{bodenY:moduli-spaces-of-parabolic-higgs-bundles}.
\begin{definition}
\source{geometric-local-systems-general-curves-isomonodromy}
\notready[Coparabolic vector bundles,
	\protect{\cite[Definition 2.3]{bodenY:moduli-spaces-of-parabolic-higgs-bundles}}]
	\label{definition:coparabolic-vector-bundle}
\source{geometric-local-systems-general-curves-isomonodromy}
	Let $E_\star$ denote a parabolic vector bundle, viewed as a parabolic
	sheaf.
	The associated {\em coparabolic vector bundle}, denoted
	$\widehat{E}_\star$, is the parabolic sheaf defined by
		$$\widehat{E}_\alpha :=
			\on{colim}_{\beta > \alpha} E_{\beta} $$
	The colimit above is the union taken over the inclusions given by
	$i^{\alpha,\beta}_E$.
\end{definition}

\begin{definition}[Coparabolic degree and slope]
\source{geometric-local-systems-general-curves-isomonodromy}
\notready
	\label{definition:}
\source{geometric-local-systems-general-curves-isomonodromy}
	If $F_\star$ is a coparabolic bundle of the form $F_\star = \widehat{E}_\star$,
for $E_\star$ a parabolic vector bundle, then the {\em coparabolic degree} of
$F_\star$ is defined by $\on{copar-deg}(F_\star) := \on{par-deg}(E_\star)$ and the
coparabolic slope of $F_\star$ is defined by $\mu_\star(F_\star) := \mu_\star(E_\star)$.
\end{definition}
\begin{example}
\source{geometric-local-systems-general-curves-isomonodromy}
\notready
	\label{example:vector-bundle-co}
\source{geometric-local-systems-general-curves-isomonodromy}
	Given a vector bundle $V$, one can define an associated parabolic bundle
	$E_\star$ with the trivial parabolic structure, as in
	\autoref{example:vector-bundle}. We call $\widehat{E}_\star$
	as the coparabolic bundle associated with trivial coparabolic structure.
\end{example}



\subsection{Induced subbundles and quotient bundles}

\label{subsubsection:induced-subbundle}
\source{geometric-local-systems-general-curves-isomonodromy}
Let $E_\star := (E, \{E^i_j\}, \{\alpha^i_j\})$ be a parabolic bundle.
	Any subbundle $F \subset E$ has an induced parabolic structure $F_\star$ as
	follows, see
\cite[Part 3, \S1.A, Definition 4]{seshadri:fibres-vectoriels-sur-les-courbes-algebriques}.
	The quasiparabolic structure over $x_j$ on $F$ is obtained
	from the filtration $$F_{x_j} = E^1_j \cap F_{x_j} \supset E^2_j \cap
	F_{x_j} \supset \cdots \supset E^{n_j+1}_j \cap F_{x_j} = 0$$
	by removing redundancies. For the weight associated to $F^i_j \subset
	F_{x_j}$ one takes $$\max_{\substack{k, 1 \leq k \leq n_j}} \{ \alpha^k_j : F^i_j = E^k_j \cap
	F_{x_j}\}.$$
	Similarly, any quotient $E \to Q$ with kernel $F$ has an induced parabolic structure
	given as follows.
	The quasiparabolic structure over $x_j$ on $Q$ is obtained from
	$$Q_{x_j} = (E^1_j+F_{x_j})/F_{x_j} \supset (E^2_j+F_{x_j})/F_{x_j}\supset\cdots\supset 
	(E^{n_j+1}_j+F_{x_j})/F_{x_j} = 0$$
	by removing redundancies. For the weight associated to a subspace $Q^i_j \subset
Q_{x_j}$ one takes $$\max_{\substack{k, 1 \leq k \leq n_j}} \{ \alpha^k_j : Q^i_j = (E^k_j + F_{x_j})
/F_{x_j}\}.$$

\subsection{Stability of parabolic and coparabolic bundles}
\label{subsubsection:stability-parabolic}
\source{geometric-local-systems-general-curves-isomonodromy}

\begin{definition}[Parabolic stability]
\source{geometric-local-systems-general-curves-isomonodromy}
\notready
	\label{definition:}
\source{geometric-local-systems-general-curves-isomonodromy}
	A parabolic vector bundle $E_\star$ is {\em parabolically semistable}
	(respectively, parabolically stable) if $\mu_\star(F_\star) \leq
\mu_\star(E_\star)$ (respectively $\mu_\star(F_\star)< \mu_\star(E_\star)$) for all parabolic subbundles $F_\star \subset E_\star$ (with the
induced parabolic structure as described above).
\end{definition}
\begin{definition}[Coparabolic stability]
\source{geometric-local-systems-general-curves-isomonodromy}
\notready
	\label{definition:coparabolic-stability}
\source{geometric-local-systems-general-curves-isomonodromy}
	A coparabolic bundle $E_\star$ is {\em coparabolically semistable}
	if $E_\star$ is of the form $\widehat{F}_\star$ for $F$ a semistable
	parabolic bundle.
\end{definition}
\begin{remark}
\source{geometric-local-systems-general-curves-isomonodromy}
\notready
	\label{remark:}
\source{geometric-local-systems-general-curves-isomonodromy}
	Note that parabolic and coparabolic stability are both defined with
	respect to parabolic bundles.
\end{remark}
\begin{remark}
\source{geometric-local-systems-general-curves-isomonodromy}
\notready
	\label{remark:coparabolic-stability}
\source{geometric-local-systems-general-curves-isomonodromy}
	This remark will not be needed in what follows.
	It turns out that if $n > 0$,
	a coparabolic bundle $E_\star=\widehat{F}_\star$ is coparabolically semistable if for any
	injection from a parabolic vector bundle $G_\star \hookrightarrow E_\star$, $\mu_\star(G_\star) <
	\mu_\star(E_\star)$.
	That is, although the definition 
	only gives $\mu_\star(G_\star) \leq \mu_\star(E_\star)$, equality of
	slopes is not possible when $n > 0$.

	The reason for this is that any such map $G_\star \to E_\star$ also
	induces a map $G[-\varepsilon]_\star \to E_\star\to F_\star$ for some sufficiently
	small $\varepsilon > 0$. 
	We then have $\mu_\star(G_\star) < \mu_\star(G[-\varepsilon]_\star) \leq
	\mu_\star(F_\star)=\mu_\star(E_\star)$.
\end{remark}



\begin{lemma}
\source{geometric-local-systems-general-curves-isomonodromy}
\notready
	\label{lemma:quotient-semistable}
\source{geometric-local-systems-general-curves-isomonodromy}
	Suppose $E_\star$ is a parabolic bundle and $F_\star \subset E_\star$ is
	a parabolic subbundle with the induced subbundle structure. If $Q_\star
	= E_\star/F_\star$ then $\on{par-deg}(E_\star) = \on{par-deg}(Q_\star) +
	\on{par-deg}(F_\star)$.
	In particular, a parabolic bundle $E_\star$ is semistable (respectively,
	stable) if and only if for every quotient
	bundle $Q_\star$, $\mu_\star(E_\star) \leq \mu_\star(Q_\star)$, (respectively $<
	\mu_\star(Q_\star)$).
\end{lemma}
\begin{proof}
	The second statement follows from the first, because
	$\frac{\on{par-deg}(F_\star)}{\rk(F_\star)} <
	\frac{\on{par-deg}(E_\star)}{\rk(E_\star)}$
	is equivalent to 
	$\frac{\on{par-deg}(E_\star)- \on{par-deg}(F_\star)}{\rk(E_\star) -
	\rk(F_\star)} >
	\frac{\on{par-deg}(E_\star)}{\rk(E_\star)}$, and similarly where one
	replaces the inequalities with equalities.

	The first statement is stated in
	\cite[Part 3, \S1.A, p. 69, Remark
	3]{seshadri:fibres-vectoriels-sur-les-courbes-algebriques},
	together with the definition of exact sequence of parabolic bundles
	\cite[Part 3, \S1.A, p. 68]{seshadri:fibres-vectoriels-sur-les-courbes-algebriques}.
% 
%	For the first statement, the equality of degrees certainly holds on the
%	underlying bundles. That is, $\deg E = \deg Q + \deg F$.
%	Write $E_\star = (E, \{E^i_j\}, \{\alpha^i_j\}),
%	F_\star = (F, \{F^i_j\}, \{\beta^i_j\}),
%	Q_\star = (Q, \{Q^i_j\}, \{\gamma^i_j\})$.
%	It suffices to show 
%	that for each $1 \leq j \leq n$, 
%	\begin{align}
%		\label{align:degree-at-point-sum}
%	\sum_i \alpha^i_j \dim (E^i_j/E^{i+1}_j) = 
%	\sum_i \beta^i_j \dim (F^i_j/F^{i+1}_j) + \sum_i \gamma^i_j
%	\dim(Q^i_j)/Q^{i+1}_j.
%	\end{align}
%	By induction on the rank of $F$, it is enough to verify this
%	equality when $\rk(F) = 1$.
%	Namely, if $S_\star \subset F_\star$ has rank $1$, we can use the claim for $F_\star/S_\star
%	\to E_\star/S_\star \to E_\star/F_\star$, $S_\star \to F_\star \to F_\star/S_\star$ and $S_\star \to E_\star \to E_\star/S_\star$ to decude the
%	claim for $F_\star \to E_\star \to E_\star/F_\star$.
%	In the case, $F$ has rank $1$, $F$ corresponds to a line $F \subset E$ indexed by a
%	slope $\beta^1_j = \alpha^s_j$, for $\alpha^s_j$ maximal among all
%	$\alpha^i_j$ for which $E^i_j$ contains the line $F$.
%	In the quotient filtration, if $\rk (E^s_j) > 1$, we have $\rk(Q^i_j) =
%	\rk(E^i_j)$ if $i \neq s$, $\rk Q^s_j = \rk E^s_j - 1$, and $\gamma^i_j
%	= \alpha^i_j$. If $\rk(E^s_j) = 1$ then $\rk Q^i_j = \rk E^i_j$ and
%	$\alpha^i_j = \gamma^i_j$ for $i <
%	s$, while $\rk Q^i_j = \rk E^{i+1}_j$ and $\gamma^i_j = \alpha^{i+1}_j$ for $i >s$.
%	In either case, \autoref{align:degree-at-point-sum} holds.
\end{proof}

\subsection{Harder-Narasimhan filtrations}
\label{subsubsection:harder-narasimhan-parabolic}
\source{geometric-local-systems-general-curves-isomonodromy}
It is a standard fact that parabolic vector bundles have a Harder-Narasimhan
filtration, and its proof is similar to the construction of Harder-Narasimhan
filtrations of vector bundles, see 
\cite[Part 3, \S1.B, Theorem 8]{seshadri:fibres-vectoriels-sur-les-courbes-algebriques}.
	
\subsection{Serre duality}

If $E_\star$ is a parabolic sheaf on a scheme
$X$, we have $H^0(C, E_\star) := \on{Hom}(\mathscr O_X, E_\star) =
\on{Hom}(\mathscr O_X, E_0) = \Gamma(X, E)$, and one can define the
higher cohomology groups by taking the corresponding right derived functor, as in 
\cite[p. 130]{yokogawa:infinitesimal-deformation}, where, more generally,
$\on{Ext}^i(E_\star, F_\star)$ is defined for parabolic $\mathscr O_X$ modules.
In general we have $H^i(X, E_\star)=H^i(X, E_0)$.

To state Serre duality, we need the notion of parabolic tensor product and
duality.
\begin{definition}
\source{geometric-local-systems-general-curves-isomonodromy}
\notready
	\label{definition:dualization}
\source{geometric-local-systems-general-curves-isomonodromy}
	For $F_\star$ a parabolic sheaf, define $F^\vee_\star :=
	\sheafhom(F_\star, \mathscr O_X)_\star$.
\end{definition}
The following gives a useful alternate description of parabolic dualization.
\begin{lemma}[Parabolic dualization, cf. \protect{\cite[(3.1)]{yokogawa:infinitesimal-deformation}}]
\source{geometric-local-systems-general-curves-isomonodromy}
\notready
	\label{lemma:parabolic-dual-formula}
\source{geometric-local-systems-general-curves-isomonodromy}
	If $F_\star$ arises from a parabolic vector bundle, we have
	$F^\vee_\alpha \simeq (\widehat{F}_{-\alpha})^\vee (-D)$.
\end{lemma}

\begin{definition}
\source{geometric-local-systems-general-curves-isomonodromy}
\notready[Parabolic tensor product, cf. \protect{\cite[Example
	3.2]{yokogawa:infinitesimal-deformation}}]
	\label{definition:parabolic-tensor}
\source{geometric-local-systems-general-curves-isomonodromy}
	Let $\tau: X - D \to X$ denote the inclusion.
	Suppose $E_\star, F_\star$ are both parabolic modules so that each $F_\alpha$
	and $E_\alpha$ is locally free,
	define $(E_\star \otimes F_\star)_\alpha := \sum_{\alpha_1 + \alpha_2 =
	\alpha} E_{\alpha_1} \otimes F_{\alpha_2}$, viewed as a subbundle of
	$\tau_* \tau^*(E \otimes F)$. We take $i^{\alpha,\beta}_{E_\star \otimes F_\star}$ to be the natural inclusion map.
\end{definition}
\begin{remark}
\source{geometric-local-systems-general-curves-isomonodromy}
\notready
	\label{remark:}
\source{geometric-local-systems-general-curves-isomonodromy}
	In \cite[p. 136-137]{yokogawa:infinitesimal-deformation}, the tensor
	product of two arbitrary parabolic sheaves is defined, but the
	definition is more difficult to state, and we will not require this
	greater generality.
\end{remark}

The next lemma states that parabolic tensor products and duals interact in the usual way with degree.
We omit the proof, which is a matter of unwinding definitions.
\begin{lemma}
\source{geometric-local-systems-general-curves-isomonodromy}
\notready
	\label{lemma:degree-in-duals}
\source{geometric-local-systems-general-curves-isomonodromy}
	For $E_\star$ and $F_\star$ two parabolic bundles, $\deg E_\star \otimes
	F_\star= \deg E_\star \rk F_\star + \deg F_\star \rk E_\star$ and $\deg
	E_\star^\vee = - \deg E_\star$.
\end{lemma}
%\begin{proof}
%	The first fact about tensor products is stated but not proved in
%	\cite[Remark, p. 5]{bodenY:moduli-spaces-of-parabolic-higgs-bundles}.
%	We content ourselves with proving 
%$\deg E_\star^\vee = - \deg E_\star$
%Using \autoref{lemma:parabolic-dual-formula}, it is enough to show the parabolic bundle 
%$F_\star$ given by 
%$F_\alpha = (\widehat{E}_{-\alpha})^\vee (-D)$ has degree $- \deg E_\star = -\deg
%E_0 - \sum_{j=1}^n \sum_{i=1}^{n_j} \alpha^i_j %\dim(E_j^i/E_j^{i+1})$.
%There are two cases, depending on whether each given $\alpha^1_j$ is $0$ or
%nonzero.
%We will cover the case that each $\alpha^1_j \neq 0$.
%The case that some $\alpha^1_j \neq 0$ is similar, but slightly more
%notationally annoying.
%One can deduce the general case by replacing $E_\star$ by a shift of the form
%$E_\star[\varepsilon]$ for which all $\alpha^1_j \neq 0$
%and using $(E_\star[\varepsilon])^\vee =
%E^\vee_\star[-\varepsilon]$ and $\deg E_\star[\varepsilon] = \deg E_\star +
%\varepsilon \rk E$.
%
%In the case all $\alpha^1_j \neq 0$, writing $F_\star = (F_0, \{F_j^i\} , \{\beta_j^i\})$,
%we find $\beta^i_j = 1 - \alpha^{n_j + 1 - i}_j$ for $1 \leq i \leq n_j$ and
%$\dim F_j^i/F_j^{i+1} = \dim E^{n_j+1-i}_j / E^{n_j + 2 -i}_j$
%for $1 \leq i \leq n_j$.
%Additionally, under our hypotheses, $F_0 = E_0^\vee(-D)$ so $\deg F_0 = -\deg
%E_0 - n\rk E$.
%Therefore,
%\begin{align*}
%	\deg F &= \deg F_0 + \sum_{j=1}^n \sum_{i=1}^{n_j} \beta^i_j
%	\dim(F_j^i/F_j^{i+1}) \\
%	&= \deg F_0 + \sum_{j=1}^n \sum_{i=1}^{n_j} (1-\alpha^{n_j+1-i}_j) \dim
%	E^{n_j+1-i}_j / E^{n_j + 2 -i}_j \\
%	&= -\deg E_0 -n \rk E+ \sum_{j=1}^n \sum_{i=1}^{n_j} (1-\alpha^{i}_j)
%	\dim E^{i}_j / E^{i + 1}_j \\
%	&= -\deg E_0 -n \rk E  + n \rk E - \sum_{j=1}^n \sum_{i=1}^{n_j}
%	(1-\alpha^{i}_j) \dim E^{i}_j / E^{i + 1}_j \\
%	&= -\deg E_0  - \sum_{j=1}^n \sum_{i=1}^{n_j} (1-\alpha^{i}_j) \dim
%	E^{i}_j / E^{i + 1}_j \\
%	&= - \deg E_\star. \qedhere
%\end{align*}
%\end{proof}



\begin{proposition}[Serre duality]
\source{geometric-local-systems-general-curves-isomonodromy}
\notready
	\label{proposition:serre-duality}
\source{geometric-local-systems-general-curves-isomonodromy}
	Suppose $X$ is a smooth projective $n$-dimensional variety over an
	algebraically closed field $k$ and let $\omega_X$ denote the dualizing
	sheaf on $X$. For all parabolic vector bundles $E_\star$, we have a canonical isomorphism
\begin{align*}
	H^i(X, E_\star) \simeq H^{n-i}(X, \widehat{E_\star^\vee} \otimes
	\omega_X(D))^\vee.
\end{align*}
\end{proposition}
\begin{proof}
	The version in \cite[Proposition
	3.7]{yokogawa:infinitesimal-deformation} states
	$\on{Ext}^i_X(E_\star, F_\star \otimes \omega_X(D)) \simeq
	\on{Ext}^{n-i}_X(F_\star, \widehat{E}_\star)^\vee.$
	Using \cite[Lemma 3.6]{yokogawa:infinitesimal-deformation}, and taking
	$E_\star = \mathscr O_X$, we find
	$H^i(X, F_\star \otimes \omega_X(D)) \simeq H^{n-i}( F_\star^\vee
		\otimes
	\widehat{\mathscr O_X})^\vee \simeq H^{n-i}(X, \widehat{F_\star^\vee})^\vee$.
	Now, taking $E_\star$ as in the statement to be 
	$F_\star \otimes \omega_X(D)$, we find $F_\star^\vee = E_\star^\vee
	\otimes \omega_X(D)$ and so
$H^i(X, E_\star) \simeq H^{n-i}(X, \reallywidehat{E_\star^\vee \otimes
\omega_X(D)})^\vee \simeq 
H^{n-i}(X, \widehat{E_\star^\vee} \otimes
\omega_X(D))^\vee.$
\end{proof}
